% GENERATED FILE. Edit canonical lessons, then run npm run generate:books.
% canonical-source-hash: fnv1a64:b8d225cf81015566
% canonical-lessons: MR-C423-kharedi-karne, MR-C423-sauda, MR-C423-kimmat, MR-C423-bhav-karne, MR-C423-vikat-ghene

\chapter{Shopping and Bargains}
\label{ch:mr-shopping-and-bargains}
\hlchaptermodality{423}

\begin{chapteropening}
\textbf{By the end of this chapter:} \emph{I can talk about shopping, deals and prices, and say I bargained and bought.}

Complete the last lesson of chapter 423: I can talk about shopping, deals and prices, and say I bargained and bought.
\end{chapteropening}

\section[\texorpdfstring{\mr{खरेदी} \mr{करणे} (kharedī karṇe)}{खरेदी करणे (kharedī karṇe)}]{\mr{खरेदी} \mr{करणे} (kharedī karṇe) — to go shopping; to purchase}

\label{lesson:MR-C423-kharedi-karne}



\begin{quote}
Before the new one: say the Marathi for bustling, then the Marathi for bustle.
\end{quote}

\subsection*{You'll want to know: \mr{खरेदी} \mr{करणे}}
\textbf{\mr{खरेदी} \mr{करणे}} — \emph{kharedī karṇe} — \textquotedblleft{}to go shopping; to purchase\textquotedblright{}.

\begin{quote}
\mr{सणासाठी} \mr{आम्ही} \mr{खूप} \mr{खरेदी} \mr{केली}.
\end{quote}

\begin{quote}
\emph{saṇāsāṭhī āmhī khūp kharedī kelī.}
\end{quote}

\begin{quote}
We did a lot of shopping for the festival.
\end{quote}

Past, with -\mr{ने} (he …): \mr{त्याने} \mr{खरेदी} \mr{केली}, \emph{tyāne kharedī kelī}.

\mr{खरेदी} is feminine, so the past is \mr{खरेदी} \mr{केली}.

\begin{cousinweb}
\textbf{\mr{खरेदी} \mr{करणे}} (\emph{kharedī karṇe}) joins \textbf{\mr{खरेदी}}, \textquotedblleft{}buying\textquotedblright{}, a loan usually traced to Persian \emph{xarīdī}, to the light verb \textbf{\mr{करणे}}, \textquotedblleft{}to do\textquotedblright{}.
\end{cousinweb}

\begin{grammarlens}[title={What you've built}]
One more word for getting around.
\end{grammarlens}

\subsection*{Your turn}
\begin{itemize}
\raggedright
  \item \emph{Say it:} \emph{kharedī karṇe}
  \item \emph{Say it:} \emph{kharedī karṇe}, once more
  \item \emph{Say it:} \emph{saṇāsāṭhī āmhī khūp kharedī kelī.} — We did a lot of shopping for the festival.
  \item \emph{Recall:} say the Marathi for bustling, then the Marathi for bustle, then say \emph{kharedī karṇe} again
\end{itemize}

\begin{tcolorbox}[breakable,colback=teal!4,colframe=teal!35!black,title={Before you move on}]
What does \mr{खरेदी} \mr{करणे} mean? (\textquotedblleft{}To go shopping; to purchase\textquotedblright{}.) Say it once more.
\end{tcolorbox}
\section[\texorpdfstring{\mr{सौदा} (saudā)}{सौदा (saudā)}]{\mr{सौदा} (saudā) — a deal; a bargain}

\label{lesson:MR-C423-sauda}



\begin{quote}
Before the new one: say the Marathi for bustle, then the Marathi for to go shopping.
\end{quote}

\subsection*{You'll want to know: \mr{सौदा}}
\textbf{\mr{सौदा}} — \emph{saudā} — \textquotedblleft{}a deal; a bargain\textquotedblright{}.

\begin{quote}
\mr{घराचा} \mr{सौदा} \mr{शेवटी} \mr{पक्का} \mr{झाला}.
\end{quote}

\begin{quote}
\emph{gharācā saudā śevṭī pakkā jhālā.}
\end{quote}

\begin{quote}
In the end the deal for the house was settled.
\end{quote}

Masculine.

\mr{सौदा} \mr{पक्का} \mr{होणे}, a deal is closed; \mr{सौदा} \mr{करणे}, to strike a deal.

\begin{cousinweb}
\textbf{\mr{सौदा}} (\emph{saudā}) is a loan usually traced to Persian \emph{saudā}, \textquotedblleft{}trade, a deal\textquotedblright{}.
\end{cousinweb}

\begin{grammarlens}[title={What you've built}]
One more word for getting around.
\end{grammarlens}

\subsection*{Your turn}
\begin{itemize}
\raggedright
  \item \emph{Say it:} \emph{saudā}
  \item \emph{Say it:} \emph{saudā}, once more
  \item \emph{Say it:} \emph{gharācā saudā śevṭī pakkā jhālā.} — In the end the deal for the house was settled.
  \item \emph{Recall:} say the Marathi for bustle, then the Marathi for to go shopping, then say \emph{saudā} again
\end{itemize}

\begin{tcolorbox}[breakable,colback=teal!4,colframe=teal!35!black,title={Before you move on}]
What does \mr{सौदा} mean? (\textquotedblleft{}A deal; a bargain\textquotedblright{}.) Say it once more.
\end{tcolorbox}
\section[\texorpdfstring{\mr{किंमत} (kimmat)}{किंमत (kimmat)}]{\mr{किंमत} (kimmat) — a price; a value}

\label{lesson:MR-C423-kimmat}



\begin{quote}
Before the new one: say the Marathi for to go shopping, then the Marathi for a deal.
\end{quote}

\subsection*{You'll want to know: \mr{किंमत}}
\textbf{\mr{किंमत}} — \emph{kimmat} — \textquotedblleft{}a price; a value\textquotedblright{}.

\begin{quote}
\mr{या} \mr{साडीची} \mr{किंमत} \mr{किती} \mr{आहे}?
\end{quote}

\begin{quote}
\emph{yā sāḍīcī kimmat kitī āhe?}
\end{quote}

\begin{quote}
What is the price of this sari?
\end{quote}

Feminine.

Also worth: \mr{त्याला} \mr{किंमत} \mr{नाही}, it has no value.

\begin{cousinweb}
\textbf{\mr{किंमत}} (\emph{kimmat}) is a loan through Persian, usually traced to Arabic \emph{qīmat}, \textquotedblleft{}value, price\textquotedblright{}.
\end{cousinweb}

\begin{grammarlens}[title={What you've built}]
One more word for getting around.
\end{grammarlens}

\subsection*{Your turn}
\begin{itemize}
\raggedright
  \item \emph{Say it:} \emph{kimmat}
  \item \emph{Say it:} \emph{kimmat}, once more
  \item \emph{Say it:} \emph{yā sāḍīcī kimmat kitī āhe?} — What is the price of this sari?
  \item \emph{Recall:} say the Marathi for to go shopping, then the Marathi for a deal, then say \emph{kimmat} again
\end{itemize}

\begin{tcolorbox}[breakable,colback=teal!4,colframe=teal!35!black,title={Before you move on}]
What does \mr{किंमत} mean? (\textquotedblleft{}A price; a value\textquotedblright{}.) Say it once more.
\end{tcolorbox}
\section[\texorpdfstring{\mr{भाव} \mr{करणे} (bhāv karṇe)}{भाव करणे (bhāv karṇe)}]{\mr{भाव} \mr{करणे} (bhāv karṇe) — to bargain}

\label{lesson:MR-C423-bhav-karne}



\begin{quote}
Before the new one: say the Marathi for a deal, then the Marathi for a price.
\end{quote}

\subsection*{You'll want to know: \mr{भाव} \mr{करणे}}
\textbf{\mr{भाव} \mr{करणे}} — \emph{bhāv karṇe} — \textquotedblleft{}to bargain\textquotedblright{}.

\begin{quote}
\mr{आजीने} \mr{भाजीवाल्याशी} \mr{भाव} \mr{केला}.
\end{quote}

\begin{quote}
\emph{ājīne bhājīvālyāśī bhāv kelā.}
\end{quote}

\begin{quote}
Grandmother bargained with the vegetable seller.
\end{quote}

Past, with -\mr{ने} (he …): \mr{त्याने} \mr{भाव} \mr{केला}, \emph{tyāne bhāv kelā}.

The seller takes -\mr{शी}: \mr{दुकानदाराशी} \mr{भाव} \mr{करणे}.

\begin{cousinweb}
\textbf{\mr{भाव} \mr{करणे}} (\emph{bhāv karṇe}) joins \textbf{\mr{भाव}}, \textquotedblleft{}the going rate\textquotedblright{}, inherited from Sanskrit \emph{bhāva}, where CDIAL lists Marathi \emph{bhāv}, \textquotedblleft{}market rate\textquotedblright{} (CDIAL 9475), to \textbf{\mr{करणे}}.
\end{cousinweb}

\begin{grammarlens}[title={What you've built}]
One more word for getting around.
\end{grammarlens}

\subsection*{Your turn}
\begin{itemize}
\raggedright
  \item \emph{Say it:} \emph{bhāv karṇe}
  \item \emph{Say it:} \emph{bhāv karṇe}, once more
  \item \emph{Say it:} \emph{ājīne bhājīvālyāśī bhāv kelā.} — Grandmother bargained with the vegetable seller.
  \item \emph{Recall:} say the Marathi for a deal, then the Marathi for a price, then say \emph{bhāv karṇe} again
\end{itemize}

\begin{tcolorbox}[breakable,colback=teal!4,colframe=teal!35!black,title={Before you move on}]
What does \mr{भाव} \mr{करणे} mean? (\textquotedblleft{}To bargain\textquotedblright{}.) Say it once more.
\end{tcolorbox}
\section[\texorpdfstring{\mr{विकत} \mr{घेणे} (vikat gheṇe)}{विकत घेणे (vikat gheṇe)}]{\mr{विकत} \mr{घेणे} (vikat gheṇe) — to buy}

\label{lesson:MR-C423-vikat-ghene}



\begin{quote}
Before the new one: say the Marathi for a price, then the Marathi for to bargain.
\end{quote}

\subsection*{You'll want to know: \mr{विकत} \mr{घेणे}}
\textbf{\mr{विकत} \mr{घेणे}} — \emph{vikat gheṇe} — \textquotedblleft{}to buy\textquotedblright{}.

\begin{quote}
\mr{मी} \mr{बाजारातून} \mr{नवीन} \mr{छत्री} \mr{विकत} \mr{घेतली}.
\end{quote}

\begin{quote}
\emph{mī bājārātūn navīn chatrī vikat ghetlī.}
\end{quote}

\begin{quote}
I bought a new umbrella at the market.
\end{quote}

Past, with -\mr{ने} (he …): \mr{त्याने} \mr{विकत} \mr{घेतलं}, \emph{tyāne vikat ghetlaṁ}.

The verb agrees with what is bought: \mr{छत्री} \mr{विकत} \mr{घेतली}, \mr{पुस्तक} \mr{विकत} \mr{घेतलं}.

\begin{cousinweb}
\textbf{\mr{विकत} \mr{घेणे}} (\emph{vikat gheṇe}) is built from \emph{vikat}, \textquotedblleft{}by purchase\textquotedblright{}, from the family of Sanskrit \emph{vikraya}, \textquotedblleft{}sale\textquotedblright{} (CDIAL 11638), and \textbf{\mr{घेणे}}, \textquotedblleft{}to take\textquotedblright{}.
\end{cousinweb}

\begin{grammarlens}[title={What you've built}]
One more word for getting around.
\end{grammarlens}

\subsection*{Your turn}
\begin{itemize}
\raggedright
  \item \emph{Say it:} \emph{vikat gheṇe}
  \item \emph{Say it:} \emph{vikat gheṇe}, once more
  \item \emph{Say it:} \emph{mī bājārātūn navīn chatrī vikat ghetlī.} — I bought a new umbrella at the market.
  \item \emph{Recall:} say the Marathi for a price, then the Marathi for to bargain, then say \emph{vikat gheṇe} again
\end{itemize}

\begin{tcolorbox}[breakable,colback=teal!4,colframe=teal!35!black,title={Before you move on}]
What does \mr{विकत} \mr{घेणे} mean? (\textquotedblleft{}To buy\textquotedblright{}.) Say it once more.
\end{tcolorbox}
